\section{Introduction}

\begin{frame}
  \frametitle{Our Work}

  \pause An implementation\footnote{\url{https://github.com/thejohncrafter/SSTT26/}}
  of a type theory in Agda.

  \pause Featuring \textbf{Definitional Proof-Irrelevance}
  \begin{tightcenter}
    Let $P$ be a proposition and $p,q:P$ be proofs. \\
    We have $p≡q$.
  \end{tightcenter}

  \pause Methods
  \begin{tightitemize}
    \item<+(1)-> A methodology for mechanizations of dependent type
      theories\myfootcite{marquet}.
    \item<+(1)-> A framework for building $n$-categories inductively from their
      $(n-1)$-categories of morphisms\myfootcite{despalungue}.
    \item<+(1)-> \textbf{Dimensional layers to control (ir)relevance.}
  \end{tightitemize}
\end{frame}

\begin{frame}
  \frametitle{This Talk}

  \pause This talk is about how thinking in terms of categorical dimension
  is useful for implementers.

  \pause I am convinced it is, because this point of view got me a simple
  solution to a difficult engineering problem.

  \begin{enumerate}
    \item<+(1)-> Reference implementations: following
      Gilbert~\textit{et.~al.}\myfootcite{gilbert}.
    \item<+(1)-> My approach to implementing type theory.
    \item<+(1)-> How the two don't match.
    \item<+(1)-> How to resolve this tension and the geometrical nature of the
      solution.
  \end{enumerate}
\end{frame}

\begin{frame}
  \frametitle{If You Only Remember One Slide}

  \pause
  low-dimensional category theory
  \begin{tightcenter}
    \begin{tabular}{ll@{}r}
      Sets         & dim. &  0 \\
      Propositions & dim. & -1
    \end{tabular}
  \end{tightcenter}

  This dimension determines the status of equality:
  
  \begin{minipage}[t]{0.4\textwidth}
    \begin{tightcenter}
    \textbf{For Propositions}
    \end{tightcenter}

    $P$ a proposition and $p,q:P$ proofs \\
    We have $p≡q$ meaning $p$ and $q$ are definitionally equal.
  \end{minipage}
  \hspace{0.09\textwidth}
  \begin{minipage}[t]{0.4\textwidth}
    \begin{tightcenter}
    \textbf{For Sets}
    \end{tightcenter}

    $X$ a set and $x,y:X$ points \\
    There is a proposition denoted $x=y$, of which proofs witness the equality
    of $x$ and $y$.
  \end{minipage}
\end{frame}

